
%%%%%%%%%%%%%%%%%%%%%%%%%%%%%%%%%%%%%%%%%%%%%%%%%%%%%%%%%%%%%%%%%%%%%%%%%%%%%%%%%%%%%%%%%%%%%%%%%%% 
% KAPITOLA 9
%%%%%%%%%%%%%%%%%%%%%%%%%%%%%%%%%%%%%%%%%%%%%%%%%%%%%%%%%%%%%%%%%%%%%%%%%%%%%%%%%%%%%%%%%%%%%%%%%%
\chapter{Limitácie programu a možné rozšírenia}
\label{ch:kap9}

Pri návrhu, implementácii a testovaní navrhnutého riešenia sa objavilo niekoľko obmedzení, ktoré sú nutné v tejto práci pomenovať. Zatiaľ čo sa Kapitola~\ref{ch:kap7} zaoberá opisom testovania a Kapitola~\ref{ch:kap8} obsahuje podrobné vyhodnotenie výsledkov, táto kapitola sa zameriava na identifikáciu limitácií nástroja (Sekcia~\ref{sub:limitacia_nastroja}), obmedzení generátora Dataster (Sekcia~\ref{sub:limitacia_generatora_dataster}) a obmedzenia pri práci s LLM (Sekcia~\ref{sub:limitacia_pri_praci_s_llm}), ako aj návrhy pre rozšírenie implementácie nástroja budúcej práci (Sekcia~\ref{sub:mozne_rozsirenia}). Dôležitou poznámkou je, že úrovne L1-L3 merajú len použiteľnosť generovaných dát a nekontrolujú sémantickú správnosť.

\section{Limitácia nástroja}
\label{sub:limitacia_nastroja}



\begin{itemize}
    \item \textbf{Zložené primárne kľúče (M:N tabuľky)}
    Tabuľky, ktorých zložený primárny kľúč je poskladaný výlučne z cudzích kľúčov (M:N), sú vylúčené už v predspracovaní so záznamom \texttt{composite\_pk}. Tabuľky so zloženým PK, ktoré toto kritérium nespĺňajú, sa spracujú, unikátnosť kombinácie kľúča však nie je garantovaná.
    
    \item \textbf{Nečíselné primárne kľúče}
    Ak primárny kľúč obsahuje textové alebo iné nečíselné hodnoty, program ich spracováva prostredníctvom zbierania vzoriek a následného generovania BLOCKS formátu. Neexistuje garancia unikátnosti, ak počet generovaných riadkov presahuje počet hodnôt v BLOCKS, vznikajú duplicitné kľúče.
    
    \item \textbf{sebareferenčná definícia v sekcii \texttt{schema} }
    Validačná vrstva nekontroluje prípad kde sa v YAML recept nachádza mapovanie stĺpca na ten istý stĺpec v rovnakej tabuľke. 

    \begin{lstlisting}
        schema:
            table_name:
                column_name: FROM table_name.column_name
    \end{lstlisting}
    
    Dôvodom je medzera v implementácii vo validačnej a opravnej vrstve, nie je to zlyhanie LLM. Tým pádom chybný recept tak prejde validáciou ako správny, ale vo vrstve L2 spôsobí pád nástroja Dataster. Tento defekt spôsobil 49 pádov pri konfigurácii qwen2.5:7b (Sekcia~\ref{sub:odpoved_rq2}). Pre opravu stačí doplnenie kontroly, ako vo validačnej, tak aj v opravnej vrstve.

    \item \textbf{Zhoršenie dátového typu pri neznámej hodnote}
    Ak nástroj vyhodnotí, že stĺpec v sekcii types má nejaký neznámy typ, alebo nebol rozpoznaný dataset z AVAILABLE\_INCLUDES, tak sa nástroj pokúsi dátový typ uhádnuť. V prípade, že sa to nepodarí, použije záložný univerzálny typ string. Toto je spôsobené nedostatočnou analýzou pri neznámych schémach. Toto vedie ku zníženiu kvality výsledného receptu, pretože dochádza ku zmene významu hodnôt vo vnútri stĺpca.

    \item \textbf{Chýbajúca používateľská definícia datasetov}
    Nástroj je limitovaný na sedem preddefinovaných datasetov v definícii AVAILABLE\_INCLUDES.  Toto je limitácia architektúry. V prípade, že sa hodnota v stĺpci v types sekcii nezhoduje ani s jedným z týchto datasetov, program zmení jeho typ na generický string. Taktiež, používateľ nemôže bez zásahu do zdrojového kódu pridať vlastné datasety. Toto zhoršuje kvalitu výsledného receptu.
    
\end{itemize}

\section{Limitácia generátora Dataster}
\label{sub:limitacia_generatora_dataster}

Nasledujúce obmedzenia boli identifikované analýzou zdrojového kódu a reprodukované na databázach testdb21-23 a extreme\_testdb6 (Kapitola~\ref{ch:kap8}). Nejde o hodnotenie nástroja ale o vymenovanie obmedzení, s ktorými musí generovanie receptov počítať.

\begin{itemize}
    \item \textbf{Escapovanie SQL hodnôt} - niektoré SQL INSERT príkazy obsahujú neescapované apostrofy a spätné lomky vo vygenerovaných reťazcoch
    \item \textbf{Zlomkové hodnoty vo vnútri finterval} - pre hodnoty s plávajúcou čiarkou sú generované hodnoty vo forme zlomkov. Toto bolo potvrdené testom \texttt{testdb21} nad DBMS MySQL a PostgreSQL.
    
    \item \textbf{SMT count} - ak recept vyžaduje viac unikátnych kombinácií, než je SMT solver schopný dodať, knižnica dbgenx vyhodí výnimku \texttt{GenerateCountError}. Avšak webové rozhranie nástroja Dataster pri identickom recepte nezlyhá úplne, ale vytvorí neúplný výstup.

    \item \textbf{RecursionError pri širokých tabuľkách} - pri extrémnych databázach, kde tabuľka obsahuje 200 stĺpcov a viac, dochádza v Datasteri k prekročeniu limitu rekurzie \texttt{RecursionError}. Toto vyvolá pád na každom recepte.

\end{itemize}


\section{Limitácie pri práci s LLM}
\label{sub:limitacia_pri_praci_s_llm}

\begin{itemize}
    \item \textbf{Kapacita menších modelov} - 
    pri vyhodnocovaní modelov sa ukázalo, že menšie modely strácajú schopnosť udržiavať štruktúru a konzistenciu YAML konfigurácie.
    
    \item \textbf{Nedeterminizmus aj pri nulovej teplote} -
    Hoci je teplota modelu nastavená na 0 a mala by generovať rovnaké výstupy, dochádza k drobným odchýlkam pri opakovaných spusteniach rovnakej úlohy.

    \item \textbf{Orezanie kontextu pri širokých schémach} -
    Pri databázových schémach s vysokým počtom stĺpcov a komplexnými väzbami narazil systém na limity kontextového okna, čo spôsobilo orezanie vstupných dát. Tým pádom nebol LLM schopný vygenerovať jednu z dôležitých sekcií.

    \item \textbf{Halucinované pomocné stĺpce} -
    Pri komplexnejších tabuľkách mal LLM tendenciu vymýšľať si neexistujúce stĺpce. Model sa pokúšal vytvoriť pomocné typy, ktoré by mohli viesť k doplneniu logiky.
\end{itemize}

\section{Možné rozšírenia}
\label{sub:mozne_rozsirenia}
Na základe spomenutých obmedzení sú navrhnuté nasledujúce rozšírenia pre budúcu prácu
\begin{itemize}
    \item rozšírenie implementácie validačnej a opravnej vrstvy o pravidlo schema spomenuté v Sekcii~\ref{sub:limitacia_nastroja}
    \item možnosť výberu cieľovej cesty a adresára pre ukladanie vygenerovaných YAML konfigurácií
    \item rozšírenie nástroja o ďalšie databázové systémy (napríklad Oracle) 
    \item priame prepojenie s existujúcou stránkou Dataster pre pohodlnejšie použitie

\end{itemize}

